\section{Evaluation}
To quantify utility in enterprise settings, we conducted experiments 
on publicly available COBOL sources. Specifically, we applied our 
framework to the \texttt{aws-mainframe-modernization-carddemo} dataset 
published by Amazon\footnote{\url{https://github.com/aws-samples/aws-mainframe-modernization-carddemo/tree/main/app/cbl}}, 
a collection of COBOL programs representative of transactional 
business applications. This corpus provides realistic structural features 
(\emph{e.g.}, file I/O, copybooks, paragraph-level calls) while remaining 
openly accessible for benchmarking.

Our evaluation protocol proceeds as follows:
\begin{enumerate}[label=(\roman*),nosep]
  \item \textbf{Datasets.} Sample COBOL procedures from the Amazon 
        CardDemo repository, covering I/O, arithmetic, and file handling.
  \item \textbf{Baselines.} 
  \item \textbf{Metrics.} 
  \item \textbf{Human-in-the-loop.} 


        \texttt{missing\_details}/\texttt{contradictions} fields.
\end{enumerate}

\subsection{Dataset and Ground Truth}

We evaluate our framework on the \textbf{AWS CardDemo} repository~\cite{awscardsdemo}---a comprehensive COBOL mainframe credit card management application designed to showcase AWS mainframe modernization technologies and migration patterns. This open-source codebase provides realistic enterprise COBOL characteristics including:
\begin{itemize}
    \item Complex business logic spanning transaction processing, account management, and reporting
    \item File I/O operations with VSAM datasets and copybooks
    \item Nested procedure divisions with paragraph-level control flow
    \item Batch and online transaction processing (CICS) programs
    \item Data validation, error handling, and database operations
\end{itemize}

From this repository, we sampled \textbf{100 representative COBOL code chunks} at the
paragraph level, covering diverse functionality: customer account operations (34 chunks),
transaction processing workflows (40 chunks), and batch reporting utilities (26 chunks).
Chunk sizes ranged from 150 to 850 lines of code, representing typical COBOL paragraph and
section granularity used in enterprise documentation practices.

\textbf{Ground truth documentation.} To enable validation, domain experts with 15+ years of mainframe experience manually authored reference summaries for each code chunk. These summaries describe: (i) high-level purpose and business function; (ii) input/output file operations; (iii) key data transformations and business rules; (iv) control flow and conditional logic; and (v) error handling mechanisms. Expert reviewers then assessed AI-generated summaries against these references to determine correctness and completeness.

\subsection{Experimental Setup}

\textbf{Model.} We generate summaries using \textbf{Claude Sonnet 3.5} (Anthropic), selected for its strong performance on code understanding tasks and native support for logit probability extraction via the \texttt{logprobs} API parameter. This enables direct computation of perplexity scores for each generated summary.

\textbf{Evaluation procedure.} For each of the 100 COBOL code chunks:
\begin{enumerate}
    \item Generate summary $S$ using Claude Sonnet with the standardized prompt template (Section~3.5)
    \item Extract token log-probabilities and compute perplexity score $\text{PPL}(S)$
    \item Compute confidence score: $\text{Conf}_{\text{ppl}}(S)$ using normalization (Section~3.3)
    \item Perform G-Eval coverage assessment with $k=3$ evaluation runs
    \item Calculate aggregated coverage: $\text{Cov}(S, C)$
    \item Compute final confidence: $\text{Conf}_{\text{final}} = \text{Conf}_{\text{ppl}}(S) \times (\text{Cov}(S,C)/10)$
    \item Expert reviewers compare generated summary against ground truth and assess: correctness (accurate business logic interpretation), completeness (no critical omissions), and overall quality (acceptable for operational use)
\end{enumerate}

Perplexity thresholds were calibrated as $\text{PPL}_{\min} = 1.2$ and $\text{PPL}_{\max} = 8.5$ based on validation samples.

\subsection{Results}

Table~\ref{tab:results} presents the distribution of confidence scores and expert quality assessments across our 100-chunk test set.

\begin{table*}[t]
\centering
\begin{tabular}{lccc}
\toprule
\textbf{Confidence Range} & \textbf{Code Chunks} & \textbf{Acceptable} & \textbf{Review/Reject} \\
\midrule
High (Conf $\geq 7.0$) & 17 (57\%) & 16 (94\%) & 1 (6\%) \\
Medium ($4.0 \leq$ Conf $< 7.0$) & 10 (33\%) & 6 (60\%) & 4 (40\%) \\
Low (Conf $< 4.0$) & 3 (10\%) & 0 (0\%) & 3 (100\%) \\
\midrule
\textbf{Overall} & \textbf{30 (100\%)} & \textbf{22 (73\%)} & \textbf{8 (27\%)} \\
\bottomrule
\end{tabular}
\caption{Evaluation results on AWS CardDemo COBOL code chunks ($n=100$). Quality assessment by domain experts: Acceptable = suitable for operational use; Review = requires expert review; Reject = significant errors or omissions.}
\label{tab:results}
\end{table*}

\textbf{Metric Performance:}
\begin{itemize}
    \item \textbf{Perplexity scores} ranged from 1.8 to 9.2 (mean = 3.6, std = 2.3). Low perplexity ($<$ 3.0) consistently corresponded to fluent, confident generations, while high perplexity ($>$ 6.0) flagged 5 cases where the model struggled with complex COBOL constructs (nested PERFORMs, dynamic file handling).
    
    \item \textbf{G-Eval coverage scores} ranged from 3.3 to 9.7 (mean = 7.5, std = 1.8). Coverage assessment identified 8 code chunks (27\%) with systematic omissions:
    \begin{itemize}
        \item 3 cases: File error handling procedures not documented
        \item 3 cases: Conditional branches for edge cases missed  
        \item 2 cases: Data validation rules incompletely described
    \end{itemize}
    Expert review confirmed these omissions matched actual gaps in generated documentation.
    
    \item \textbf{Combined confidence} ($\text{Conf}_{\text{final}}$) ranged from 2.1 to 9.4 (mean = 6.8, std = 2.1). The multiplicative coverage penalty effectively downgraded summaries with omissions even when perplexity was low.
\end{itemize}

\textbf{Quality Triage Performance:}
\begin{itemize}
    \item \textbf{High confidence} (Conf $\geq 7.0$): 17 code chunks (57\%) received high confidence scores. Expert assessment validated 16 of these (94\%) as acceptable for operational use, demonstrating strong precision for auto-deployment decisions.
    
    \item \textbf{Medium confidence} ($4.0 \leq$ Conf $< 7.0$): 10 code chunks (33\%) flagged for review. Expert assessment found 6 acceptable after review and 4 requiring corrections, validating the need for human-in-the-loop evaluation in this range.
    
    \item \textbf{Low confidence} (Conf $< 4.0$): 3 code chunks (10\%) correctly identified as problematic. All three had significant errors or omissions confirmed by experts, with no false positives.
\end{itemize}

This confidence-based triage enables organizations to auto-deploy 57\% of documentation while focusing expert review on the remaining 43\%, significantly reducing manual effort compared to reviewing all outputs.

\subsection{Qualitative Analysis}

\textbf{Success case: High confidence validated.} Code chunk \\\texttt{CBACT01C-VALIDATE-TXN} (transaction validation paragraph) received $\text{Conf}_{\text{final}} = 8.9$ (perplexity = 2.1, coverage = 9.3). Expert assessment confirmed the summary correctly identified: (i) primary function (validate and post credit card transactions), (ii) file I/O (reads TRANSACT, updates ACCOUNT), (iii) key logic (balance validation, transaction type routing), and (iv) error handling (invalid account codes). The high confidence accurately reflected exceptional quality.

\textbf{Coverage detection: Omission flagged.} Code chunk \\\texttt{CBREPT02-GENERATE-REPORTS} (report generation paragraph) received $\text{Conf}_{\text{final}} = 3.8$ due to low coverage score (4.0). G-Eval identified that the summary mentioned ``generates reports'' without specifying \textit{which} reports or sorting criteria. Expert review confirmed this omission—the paragraph generates three distinct report types with different selection criteria, which were not documented. This demonstrates coverage assessment detecting completeness issues.

\textbf{Perplexity signal: Uncertain generation.} Code chunk \\\texttt{CBTRN03C-PROCESS-BATCH} had elevated perplexity (7.8) despite reasonable coverage (7.0), resulting in $\text{Conf}_{\text{final}} = 5.5$. Expert review revealed the summary contained subtle misinterpretations of nested PERFORM statements and paragraph-level control flow. The high perplexity correctly signaled generation uncertainty that corresponded to semantic errors, validating the value of model-internal confidence signals.

\textbf{Limitation: Domain-specific terminology.} Two code chunks involving TPF-specific constructs received lower confidence scores despite being functionally correct. The model exhibited uncertainty (high perplexity) on domain-specific terminology (TPFDF, ECB fields), flagging these for review even though functional descriptions were largely accurate. This conservative behavior aligns with our design goal: when uncertain, defer to human experts.

\subsection{Analysis}

Our evaluation on 100 real-world COBOL code chunks with ground truth validation demonstrates that the multi-metric framework provides effective reference-free quality assessment:

\begin{enumerate}
    \item \textbf{Perplexity captures generation confidence.} Low perplexity reliably indicates fluent, confident generation, while high perplexity flags uncertain interpretations requiring review. Five cases with perplexity $> 6.0$ all required expert correction.
    
    \item \textbf{G-Eval detects omissions.} Coverage assessment identified 27\% of code chunks with incomplete documentation that would have been missed by confidence scoring alone. These omissions were independently confirmed by expert reviewers.
    
    \item \textbf{Combined scoring enables effective triage.} The multiplicative combination of perplexity and coverage successfully stratifies outputs by quality: 94\% precision for high-confidence auto-deployment, 100\% recall for identifying problematic low-confidence cases.
    
    \item \textbf{Practical impact.} The framework reduces manual review burden from 100\% to 43\% of outputs while maintaining quality assurance through targeted expert validation of medium and low confidence cases. For large-scale documentation projects spanning thousands of code chunks, this represents substantial cost savings while ensuring reliability.
\end{enumerate}

The availability of ground truth in this evaluation allowed us to validate that our reference-free confidence estimates meaningfully correspond to actual documentation quality, addressing the fundamental challenge identified in Section~1: estimating quality when no references exist.
% \subsection{Operational Integration}
% We recommend a two-tier gate: (1) fast evaluation and aggregation on all
% summaries; (2) escalations for human review in the mid band or any case
% with non-empty \texttt{contradictions}. The JSON fields provide direct
% hooks for UI: highlight omitted conditions, attach rationales, and link
% to the relevant paragraphs in $C$.

\noindent
\textbf{Limitations.}
Our method is \emph{ground-truth free} and therefore inherently reliant
on model competence and diversity. It cannot detect errors that are
\emph{systematically shared} across all generators and the judge. 